% 第4章 针尖纳米结构的手性表面波动力学
\chapter{针尖纳米结构的手性表面波动力学}
\label{chap:result1}

\section{引言}
\label{sec:r1-intro}

% 纳米针尖作为表面等离激元研究的经典模型体系
% 手性光-物质相互作用的基本概念
% 本章研究目标：利用阿秒电子显微术观测针尖上手性表面波的亚周期动力学

\section{纳米针尖的制备与表征}
\label{sec:r1-fabrication}

% 钨针尖的FIB加工
% 狭缝结构的引入
% SEM表征

\begin{figure}[htbp]
    \centering
    % 占位符：用户实验数据图
    \fbox{\parbox{0.8\linewidth}{\centering\vspace{3cm} 图4-1：钨纳米针尖SEM图像\\（实验数据待插入）\vspace{3cm}}}
    \caption{\label{fig:ch4-tip} 狭缝修饰的钨纳米针尖。(a) 全景SEM图像；(b) 针尖顶端狭缝结构的高倍率图像。}
\end{figure}

\section{圆偏振光激发的手性响应}
\label{sec:r1-chiral}

% 左旋和右旋圆偏振光分别激发
% 方向性表面波的产生
% 手性选择性机制：自旋-轨道耦合

\section{亚周期时间分辨测量}
\label{sec:r1-temporal}

% 时间延迟扫描数据
% 表面波传播方向的手性控制
% 相位速度测量

\begin{figure}[htbp]
    \centering
    % 占位符：用户实验数据图
    \fbox{\parbox{0.8\linewidth}{\centering\vspace{3cm} 图4-2：手性表面波时间分辨测量\\（实验数据待插入）\vspace{3cm}}}
    \caption{\label{fig:ch4-chiral} 圆偏振光激发下手性表面波的时间分辨测量。(a-b) 左旋（LCP）和右旋（RCP）圆偏振光激发下的能量滤波图像；(c-e) 选定延迟时刻的场演化。}
\end{figure}

\section{表面波传播特性的定量分析}
\label{sec:r1-analysis}

% 相位速度的超光速特征
% 近场衰减距离
% 左旋-右旋之间$\sim$1.8\,fs的延迟

\section{本章小结}
\label{sec:r1-summary}

本章利用阿秒电子显微术研究了狭缝修饰钨针尖上的手性表面波动力学。实验结果表明，圆偏振光的偏振手性可以通过自旋-轨道耦合机制控制表面波的传播方向，在亚周期时间尺度上实现了对纳米光子激发的方向性操控。左旋和右旋激发图案之间存在$\sim$\SI{1.8}{fs}的时间延迟，反映了手性表面波的矢量特征。
